\subsection{p 次幂可积函数的性质}

定理 4. --- 设 F 与 G 是两个 Banach 空间，u 是 F 到 G 中的一个连续线性映射。对每个函数 \(f \in L_{F}^{p}\) ，复合函数 \(u \circ f\) 属于 \(L_{G}^{p}\) 。

设 \(\mathbf{f} \in \mathcal{L}_{\mathrm{F}}^{p}\) ；对每个 \(\varepsilon > 0\) ，存在函数 \(\mathbf{g} \in \mathcal{K}_{\mathrm{F}}\) 使 \(\mathrm{N}_p(\mathbf{f} - \mathbf{g}) \leqslant \varepsilon\) ；由于 \(|\mathbf{u} \circ \mathbf{f} - \mathbf{u} \circ \mathbf{g}| \leqslant \| \mathbf{u} \| \cdot |\mathbf{f} - \mathbf{g}|\) ，我们有

\[
\mathrm{N} _ {p} (\mathbf {u} \circ \mathbf {f} - \mathbf {u} \circ \mathbf {g}) \leqslant \| \mathbf {u} \| \cdot \mathrm{N} _ {p} (\mathbf {f} - \mathbf {g}) \leqslant \varepsilon \| \mathbf {u} \|,
\]

而由于 \(u \circ g\) 是具有紧支集的连续函数，定理得证。

推论 1. --- 设 \(\mathbf{a}'\) 是 \(\mathbf{F}\) 上的一个连续线性形式；若 \(\mathbf{f} \in \mathcal{L}_{\mathrm{F}}^{p}\) ，则数值函数 \(x \mapsto \langle \mathbf{f}(x), \mathbf{a}' \rangle\) （记作 \(\langle \mathbf{f}, \mathbf{a}' \rangle\) ）属于 \(\mathcal{L}^p\) 。

推论 2. --- 给定 \(n\) 个点 \(\mathbf{a}_k\) （\(F (1 \leqslant k \leqslant n)\) 中的）以及 \(n\) 个数值函数 \(f_k (1 \leqslant k \leqslant n)\) （属于 \(\mathcal{L}^p\) ），则函数 \(\mathbf{f} = \sum_{k=1}^{n} \mathbf{a}_k f_k\) 属于 \(\mathcal{L}_F^p\) 。

这由 R 到 F 中的映射 \(t \mapsto at\) 是连续的这一事实得出。

命题 9. --- 设 F 是 R 上的一个 n 维向量空间，且设 \((\mathbf{e}_{k})_{1\leqslant k\leqslant n}\) 是 F 的一个基。为使函数 \(f=\sum_{k=1}^{n}\mathbf{e}_{k}f_{k}\) 属于 \(L_{F}^{p}\) ，必须且只须每个数值函数 \(f_{k}\) 都属于 \(L^{p}\) 。

这由定理 4 的推论 1 与推论 2 立得。

命题 10. --- 在空间 \(\mathcal{L}_{\mathrm{F}}^{p}\) 中，由（有限）线性组合 \(\sum_{k} \mathbf{a}_{k} f_{k}\) （其中 \(\mathbf{a}_{k} \in \mathbf{F}\) 而 \(f_{k}\) 是具有紧支集的连续数值函数）所构成的线性子空间是稠密的（对于 \(p\) 阶平均收敛拓扑）。

具有紧支集的 X 到 F 中的连续映射所成之集 \(H_{F}\) 依定义在 \(L_{F}^{p}\) 中稠密。另一方面，每个函数 \(g \in H_{F}\) 都可以用形如 \(\sum_{k} a_{k} f_{k}\) 的函数一致逼近，其中 \(f_{k}\) 是支集含于 g 的支集的一个固定紧邻域中的连续函数（Ch. III, §1, No.~2, 引理 2）；由此（No.~3, 命题 4）g 处于 \(L_{F}^{p}\) 中 \(\sum_{k} a_{k} f_{k}\) 所成之集的闭包之内，命题得证。

命题 11. --- 若函数 f 属于 \(L_{F}^{p}\) ，则函数 \(|f|\) 属于 \(L^{p}\) ，且映射 \(f \mapsto |f|\) （\(L_{F}^{p}\) 到 \(L^{p}\) 中的）是一致连续的（对于 p 阶平均收敛拓扑）。

对每个 \(\varepsilon > 0\) ，存在具有紧支集的连续函数 \(\mathbf{g}\) ，使 \(\mathrm{N}_p(\mathbf{f} - \mathbf{g}) \leqslant \varepsilon\) ；由于 \(||\mathbf{f}| - |\mathbf{g}|| \leqslant |\mathbf{f} - \mathbf{g}|\) ，我们有 \(\mathrm{N}_p(|\mathbf{f}| - |\mathbf{g}|) \leqslant \varepsilon\) ，这就证明了 \(|\mathbf{f}| \in \mathcal{L}^p\) 。另一方面，若 \(\mathbf{f}_1, \mathbf{f}_2\) 是 \(\mathcal{L}_{\mathrm{F}}^p\) 中的两个函数，则 \(\mathrm{N}_p(|\mathbf{f}_1| - |\mathbf{f}_2|) \leqslant \mathrm{N}_p(\mathbf{f}_1 - \mathbf{f}_2)\) ，这表明 \(\mathbf{f} \mapsto |\mathbf{f}|\) 是一个一致连续映射。

命题 12. --- 为使数值函数 f 属于 \(L^{p}\) ，必须且只须函数 \(f^{+}\) 与 \(f^{-}\) 的每一个都属于 \(L^{p}\) 。

条件是充分的，因为 \(f = f^{+} - f^{-}\) ；它是必要的，因为若 \(f \in \mathcal{L}^p\) 则 \(|f| \in \mathcal{L}^p\) （命题 11）。

推论. --- \(\mathcal{L}^{p}\) 中函数的有限族的上包络（相应地，下包络）属于 \(\mathcal{L}^{p}\) 。

